\documentclass[11pt]{article}
\usepackage[utf8]{inputenc}
\usepackage[T1]{fontenc}
\usepackage{multicol}
\usepackage{geometry}
\usepackage{enumitem}
\usepackage{xcolor}
\usepackage{titlesec}
\usepackage{ragged2e}

\geometry{a4paper, left=1cm, right=1cm, top=0.8cm, bottom=1.2cm}

\setlength{\parindent}{0pt}
\setlength{\columnsep}{1cm}
\setlength{\columnseprule}{0.4pt}

\titleformat{\section}{\normalfont\Large\bfseries\color{blue}}{}{0em}{}

\title{\textbf{\textcolor{red}{Funções de Primeiro Grau}}}
\author{Professor: Jefferson}
\date{}

\begin{document}

\maketitle
\vspace{-0.5cm}

\begin{justify}
\textbf{Observação: \\ Copiar e responder no caderno. Resolva os problemas utilizando a definição de função de primeiro grau, que é uma função onde o coeficiente linear é constante.}
\end{justify}

\begin{multicols}{2}

\section*{Exercícios}

\begin{enumerate}[leftmargin=*, itemsep=0pt]

\item Se $f(x) = 3x + 2$, determine $f(4)$.

\item Encontre o valor de $x$ para que $f(x) = 17$, onde $f(x) = 5x - 10$.

\item Uma função é definida por $g(x) = 4x - 3$. Calcule $g(-2)$.

\item Se $h(x) = x^2 + 6x + 9$, determine $h(0)$.

\item Encontre o valor de $x$ para que $h(x) = 15$, onde $h(x) = 2x^2 - 3x + 7$.

\item Uma função é definida por $k(x) = \frac{1}{2}x + 5$. Calcule $k(6)$.

\item Se $m(x) = x - 4$, determine $m(-1)$.

\item Encontre o valor de $x$ para que $m(x) = 0$, onde $m(x) = -x + 3$.

\item Uma função é definida por $n(x) = 3x^2 - 6x$. Calcule $n(3)$.

\item Se $p(x) = 2x + 1$, determine $p(-5)$.

\item Encontre o valor de $x$ para que $p(x) = -7$, onde $p(x) = x + 4$.

\end{enumerate}

\end{multicols}

\end{document}

